\section{Methods}

This investigation follows a structured methodology designed to minimise bias and maximise the accuracy and
reproducibility of results.
Synthetic data are generated using a custom C\texttt{++} simulator across a variable number of timesteps and
procedurally constructed initial problem configurations.

\subsection{Procedural Generation}

Initial conditions are procedurally generated using a Mersenne--Twister pseudo-random number engine with fixed seed
parameters to ensure reproducibility.
Two physical quantities require generation for each system: the mass of the central object and the spatial coordinates
of the orbiting body.

\paragraph{Central body mass.}
The stellar mass is sampled uniformly from the range $[1.0,\, 2.4]\,M_\odot$.
This interval is chosen to restrict the simulation domain to main-sequence stars of spectral classes G through A,
encompassing objects analogous to the Sun and Vega, respectively.
This restriction is physically motivated: the gravitational influence of more massive stars (class O or B) would cause
rapid orbital decay or ejection, while sub-solar-mass stars would introduce qualitatively different dynamical
regimes outside the scope of this study.

\paragraph{Planetary mass.}
The mass of each orbiting body is derived from the stellar mass via the fixed ratio
\begin{equation}
    m_\mathrm{planet} = \frac{M_\star}{333{,}000},
\end{equation}
which approximates the observed mass ratio between Earth and the Sun, providing a physically realistic planetary-mass
analogue for each generated system.

\paragraph{Orbital coordinates.}
The initial separation between the central body and the orbiting object is sampled uniformly from the
range $[1,\, 5]\,\mathrm{AU}$.
This interval is deliberately chosen to prevent premature convergence in the numerical integrator and to ensure a
well-conditioned, stable dataset throughout the integration period.

\subsection{Simulation}

Following procedural generation, the derived initial quantities velocity, acceleration, and separation distance
are computed analytically, alongside baseline stability diagnostics for angular momentum and total energy.

Simulations are executed in parallel across six CPU cores (Apple Silicon M2) with one independent simulation per core.
Two distinct integration regimes are employed, controlled by both the timestep $\Delta t$ and the total number of steps $N$:

\begin{itemize}
    \item \textbf{Coarse regime:} $N = 10^5$ steps with $\Delta t = 10^{-2}$, providing a computationally efficient initial pass.
    \item \textbf{Fine regime:} $N = 10^6$ steps with $\Delta t = 10^{-3}$, applied to an independently generated problem configuration to improve resolution and expose long-term integrator behaviour.
\end{itemize}

The reduction in $\Delta t$ for the finer regime is deliberate: maintaining the geometric structure of the integration (i.e.\ symplecticity) requires that the product $N \cdot \Delta t$ remains consistent across regimes. This dual-resolution strategy improves the fidelity of the dataset while revealing convergence limitations inherent to specific integrators.

All simulation output is written to individual \texttt{.csv} files.
Post-processing and visualisation are performed using \texttt{matplotlib} under Python~3.12.
The following plots are produced for each integrator:
\begin{enumerate}
    \item Orbital trajectory in the $xy$-plane.
    \item Angular momentum as a function of time.
    \item Total energy as a function of time.
    \item Relative error in both conserved quantities, aggregated across integrators and plotted against timestep size.
\end{enumerate}

\subsection{Integrator Derivation}

In line with one of the primary objectives of this research the development and evaluation of novel numerical schemes,
the best-performing integrator identified among the five candidates will be extended to a higher order via appropriate
composition or correction methods.
This derived integrator will then be benchmarked against its lower-order counterpart under identical simulation conditions.

\subsection{Hypothesis}

Drawing on prior analytical work and the existing literature on geometric numerical integration, the primary hypothesis
is that the \textit{Velocity Verlet} integrator will deliver the most favourable trade-off between computational
efficiency and energy fidelity.
This expectation is grounded in two properties of the method: its symplecticity, which guarantees long-term conservation
of the symplectic structure of Hamiltonian systems, and its second-order accuracy, which provides a sufficient
resolution for orbital dynamics without the overhead of higher-order schemes.

Integrator performance will be evaluated against three criteria: (i) long-term conservation of total energy, (ii)
preservation of angular momentum, and (iii) computational cost, measured in wall-clock time per integration step.